\section{Exercices d'application}

\rubrique{Calcul du produit scalaire}

\exo{}\diff{1}
Dans un repère orthonormé $(O;\vec{i},\vec{j})$, on donne $\vec{u}(2; -3)$ et $\vec{v}(-1; 4)$.
Calculer $\vec{u}\cdot\vec{v}$.

\exo{}\diff{1}
Soient $\vec{u}$ et $\vec{v}$ deux vecteurs tels que $\|\vec{u}\|=3$, $\|\vec{v}\|=5$ et $\vec{u}\cdot\vec{v}=7$.
Calculer $(\vec{u}+\vec{v})\cdot(\vec{u}-\vec{v})$.

\exo{}\diff{1}
Dans un triangle $ABC$ rectangle en $A$, on a $AB=4$ et $AC=3$. Calculer $\overrightarrow{AB}\cdot\overrightarrow{AC}$ et $\overrightarrow{BA}\cdot\overrightarrow{BC}$.

\exo{}\diff{2}
On considère un triangle $ABC$ tel que $AB=6$, $AC=8$ et $\widehat{BAC}=60^\circ$.
Calculer $\overrightarrow{AB}\cdot\overrightarrow{AC}$ et $BC$.

\rubrique{Orthogonalité}

\exo{}\diff{1}
Dans un repère orthonormé, on donne $\vec{u}(3; m)$ et $\vec{v}(2; -6)$. Déterminer $m$ pour que $\vec{u}\perp\vec{v}$.

\exo{}\diff{2}
Soit $ABC$ un triangle avec $A(1;2)$, $B(4;3)$ et $C(2;5)$.
Montrer que $ABC$ est rectangle en $A$.

\exo{}\diff{2}
Déterminer une équation de la droite passant par $A(-2;3)$ et perpendiculaire à la droite $(D): 3x-4y+5=0$.

\rubrique{Relations métriques dans le triangle}

\exo{}\diff{2}
Dans un triangle $ABC$, on a $AB=7$, $AC=9$ et $BC=8$.
Calculer $\cos\widehat{A}$ et l'angle $\widehat{A}$ (au degré près).

\exo{}\diff{2}
Soit $ABC$ un triangle tel que $AB=5$, $AC=6$ et $\widehat{A}=120^\circ$.
Calculer $BC$ et l'aire du triangle.

\exo{}\diff{3}
Dans un triangle $ABC$, on donne $AB=4$, $AC=6$ et $BC=7$.
Calculer la longueur de la médiane issue de $A$.

\rubrique{Équations de droites et de cercles}

\exo{}\diff{2}
Déterminer l'équation du cercle de centre $\Omega(2;-3)$ et de rayon $5$.

\exo{}\diff{2}
On considère les points $A(1;2)$ et $B(5;4)$. Déterminer l'équation du cercle de diamètre $[AB]$.

\exo{}\diff{3}
Soit $C$ l'ensemble des points $M(x;y)$ tels que $x^2+y^2-4x+6y-12=0$.
Montrer que $C$ est un cercle dont on précisera le centre et le rayon.

\rubrique{Lignes de niveau}

\exo{}\diff{2}
Soient $A$ et $B$ deux points tels que $AB=4$.
Déterminer l'ensemble des points $M$ du plan tels que $\overrightarrow{MA}\cdot\overrightarrow{MB}=0$.

\exo{}\diff{3}
On donne $A(1;2)$ et $B(5;2)$. Déterminer l'ensemble des points $M$ tels que $\overrightarrow{MA}\cdot\overrightarrow{MB}=4$.

\exo{}\diff{3}
Soit $ABC$ un triangle équilatéral de côté $a$.
Déterminer l'ensemble des points $M$ tels que $\overrightarrow{MA}\cdot\overrightarrow{MB}= \overrightarrow{MA}\cdot\overrightarrow{MC}$.

\section{Exercices d'entraînement}

\rubrique{Calcul du produit scalaire}

\exo{}\diff{1}
Calculer $\vec{u}\cdot\vec{v}$ dans les cas suivants :
\begin{enumerate}[label=\alph*)]
\item $\|\vec{u}\|=4$, $\|\vec{v}\|=6$, $(\vec{u},\vec{v})=45^\circ$.
\item $\vec{u}(-2;5)$, $\vec{v}(3;-1)$.
\item $\vec{u}=2\vec{i}-3\vec{j}$, $\vec{v}=-\vec{i}+4\vec{j}$.
\end{enumerate}

\exo{}\diff{1}
Soient $\vec{u}$ et $\vec{v}$ tels que $\|\vec{u}\|=2$, $\|\vec{v}\|=3$ et $\vec{u}\cdot\vec{v}=4$.
Calculer $\|\vec{u}+\vec{v}\|$ et $\|\vec{u}-\vec{v}\|$.

\exo{}\diff{2}
Dans un carré $ABCD$ de côté $4$, calculer $\overrightarrow{AB}\cdot\overrightarrow{AC}$, $\overrightarrow{AB}\cdot\overrightarrow{AD}$ et $\overrightarrow{AC}\cdot\overrightarrow{BD}$.

\exo{}\diff{2}
Soit $ABC$ un triangle équilatéral de côté $a$. Calculer $\overrightarrow{AB}\cdot\overrightarrow{AC}$ et $\overrightarrow{AB}\cdot\overrightarrow{BC}$.

\exo{}\diff{2}
Dans un repère orthonormé, on donne $\vec{u}(2;1)$, $\vec{v}(-1;3)$ et $\vec{w}(4;-2)$.
Calculer $(\vec{u}+\vec{v})\cdot\vec{w}$ et $\vec{u}\cdot(\vec{v}-\vec{w})$.

\rubrique{Orthogonalité}

\exo{}\diff{1}
Déterminer $m$ pour que $\vec{u}(m;2)$ et $\vec{v}(3;m)$ soient orthogonaux.

\exo{}\diff{2}
Soient $A(2;1)$, $B(5;3)$ et $C(4;6)$. Montrer que $ABC$ est un triangle rectangle.

\exo{}\diff{2}
Déterminer une équation de la droite passant par $A(3;-1)$ et perpendiculaire à $(D): 2x+5y-3=0$.

\exo{}\diff{2}
Soit $(D): 3x-2y+4=0$. Déterminer l'équation de la droite passant par $O$ et perpendiculaire à $(D)$.

\exo{}\diff{3}
On considère les points $A(1;2)$, $B(3;4)$ et $C(5;0)$. Déterminer l'orthocentre du triangle $ABC$.

\rubrique{Relations métriques dans le triangle}

\exo{}\diff{2}
Dans un triangle $ABC$, $AB=8$, $AC=10$ et $\widehat{A}=30^\circ$. Calculer $BC$.

\exo{}\diff{2}
Soit $ABC$ un triangle avec $AB=6$, $AC=9$ et $BC=11$. Calculer $\cos\widehat{A}$ et $\widehat{A}$.

\exo{}\diff{2}
Dans un triangle $ABC$, $AB=5$, $BC=7$ et $\widehat{B}=60^\circ$. Calculer $AC$.

\exo{}\diff{3}
Soit $ABC$ un triangle tel que $AB=3$, $AC=4$ et $BC=5$. Calculer les longueurs des trois médianes.

\exo{}\diff{3}
Dans un triangle $ABC$, $AB=7$, $AC=8$ et $BC=9$. Calculer l'aire du triangle.

\rubrique{Équations de droites et de cercles}

\exo{}\diff{2}
Déterminer l'équation du cercle de centre $\Omega(-1;4)$ et passant par $A(2;1)$.

\exo{}\diff{2}
On donne $A(2;3)$ et $B(6;5)$. Déterminer l'équation du cercle de diamètre $[AB]$.

\exo{}\diff{2}
Montrer que $x^2+y^2-6x+4y+9=0$ est l'équation d'un cercle. Préciser son centre et son rayon.

\exo{}\diff{3}
Déterminer l'équation du cercle passant par $A(1;2)$, $B(3;4)$ et $C(5;2)$.

\exo{}\diff{3}
Soit $(D): x-2y+3=0$. Déterminer l'équation du cercle de centre $\Omega(2;1)$ tangent à $(D)$.

\rubrique{Lignes de niveau}

\exo{}\diff{2}
Soient $A$ et $B$ deux points tels que $AB=6$. Déterminer l'ensemble des points $M$ tels que $\overrightarrow{MA}\cdot\overrightarrow{MB}=0$.

\exo{}\diff{2}
On donne $A(1;1)$ et $B(5;1)$. Déterminer l'ensemble des points $M$ tels que $\overrightarrow{MA}\cdot\overrightarrow{MB}=3$.

\exo{}\diff{3}
Soit $ABC$ un triangle rectangle en $A$ avec $AB=3$, $AC=4$. Déterminer l'ensemble des points $M$ tels que $\overrightarrow{MA}\cdot\overrightarrow{MB}= \overrightarrow{MA}\cdot\overrightarrow{MC}$.

\exo{}\diff{3}
On donne $A(2;3)$ et $B(6;3)$. Déterminer l'ensemble des points $M$ tels que $MA^2-MB^2=8$.

\exo{}\diff{3}
Soit $ABC$ un triangle équilatéral de côté $a$. Déterminer l'ensemble des points $M$ tels que $\overrightarrow{MA}\cdot\overrightarrow{MB}= \overrightarrow{MB}\cdot\overrightarrow{MC}$.

\section{Problèmes type Probatoire/Bac}

\sujetbac[Probatoire Blanc — Produit scalaire et géométrie]
Dans un repère orthonormé $(O;\vec{i},\vec{j})$, on considère les points $A(2;1)$, $B(5;3)$ et $C(3;5)$.
\begin{enumerate}[label=\arabic*)]
\item Calculer $\overrightarrow{AB}\cdot\overrightarrow{AC}$. En déduire la nature du triangle $ABC$.
\item Déterminer une équation du cercle $(C)$ de diamètre $[BC]$.
\item Soit $(D)$ la droite passant par $A$ et perpendiculaire à $(BC)$. Déterminer une équation de $(D)$.
\item Déterminer les coordonnées du point $H$, intersection de $(D)$ et de $(BC)$. Que représente $H$ pour le triangle $ABC$ ?
\item Calculer l'aire du triangle $ABC$.
\end{enumerate}

\sujetbac[Baccalauréat Blanc — Lignes de niveau]
On considère un triangle $ABC$ rectangle en $A$ avec $AB=4$ et $AC=3$.
\begin{enumerate}[label=\arabic*)]
\item Calculer $\overrightarrow{AB}\cdot\overrightarrow{AC}$ et $BC$.
\item Soit $I$ le milieu de $[BC]$. Montrer que pour tout point $M$ du plan, $MB^2+MC^2=2MI^2+\dfrac{BC^2}{2}$.
\item Déterminer l'ensemble $(E_1)$ des points $M$ tels que $MB^2+MC^2=25$.
\item Déterminer l'ensemble $(E_2)$ des points $M$ tels que $\overrightarrow{MA}\cdot\overrightarrow{MB}=0$.
\item Déterminer l'ensemble $(E_3)$ des points $M$ tels que $\overrightarrow{MA}\cdot\overrightarrow{MB}= \overrightarrow{MA}\cdot\overrightarrow{MC}$.
\end{enumerate}

\sujetbac[Probatoire — Relations métriques]
Dans un triangle $ABC$, on donne $AB=6$, $AC=8$ et $\widehat{A}=60^\circ$.
\begin{enumerate}[label=\arabic*)]
\item Calculer $BC$.
\item Calculer $\cos\widehat{B}$ et $\cos\widehat{C}$.
\item Calculer la longueur de la médiane issue de $A$.
\item Calculer l'aire du triangle $ABC$.
\item Soit $H$ le projeté orthogonal de $B$ sur $(AC)$. Calculer $AH$ et $BH$.
\end{enumerate}

\sujetbac[Baccalauréat — Cercle et droite]
Dans un repère orthonormé $(O;\vec{i},\vec{j})$, on considère les points $A(1;2)$ et $B(5;4)$.
\begin{enumerate}[label=\arabic*)]
\item Déterminer l'équation du cercle $(C)$ de diamètre $[AB]$.
\item Vérifier que le point $C(3;5)$ appartient à $(C)$.
\item Déterminer une équation de la tangente $(T)$ à $(C)$ au point $C$.
\item Soit $(D): 2x-y+3=0$. Montrer que $(D)$ est tangente à $(C)$.
\item Déterminer les coordonnées du point de contact.
\end{enumerate}

\section{Corrigés}

\corrige{de l'exercice 1}
$\vec{u}\cdot\vec{v}=2\times(-1)+(-3)\times4=-2-12=-14$.

\corrige{de l'exercice 2}
$(\vec{u}+\vec{v})\cdot(\vec{u}-\vec{v})=\vec{u}\cdot\vec{u}-\vec{u}\cdot\vec{v}+\vec{v}\cdot\vec{u}-\vec{v}\cdot\vec{v}=\|\vec{u}\|^2-\|\vec{v}\|^2=9-25=-16$.

\corrige{de l'exercice 3}
Dans un triangle rectangle en $A$, $\overrightarrow{AB}\cdot\overrightarrow{AC}=0$ car les vecteurs sont orthogonaux.
$\overrightarrow{BA}\cdot\overrightarrow{BC}=BA\times BC\times\cos\widehat{B}$.
$BC=\sqrt{4^2+3^2}=5$, $\cos\widehat{B}=\dfrac{AB}{BC}=\dfrac{4}{5}$.
Donc $\overrightarrow{BA}\cdot\overrightarrow{BC}=4\times5\times\dfrac{4}{5}=16$.

\corrige{de l'exercice 4}
$\overrightarrow{AB}\cdot\overrightarrow{AC}=AB\times AC\times\cos60^\circ=6\times8\times\dfrac12=24$.
D'après la formule d'Al-Kashi : $BC^2=AB^2+AC^2-2AB\cdot AC\cos\widehat{A}=36+64-2\times6\times8\times\dfrac12=100-48=52$.
Donc $BC=\sqrt{52}=2\sqrt{13}$.

\corrige{de l'exercice 5}
$\vec{u}\perp\vec{v}\iff\vec{u}\cdot\vec{v}=0\iff3\times2+m\times(-6)=0\iff6-6m=0\iff m=1$.

\corrige{de l'exercice 6}
$\overrightarrow{AB}(3;1)$, $\overrightarrow{AC}(1;3)$.
$\overrightarrow{AB}\cdot\overrightarrow{AC}=3\times1+1\times3=6\neq0$, donc $ABC$ n'est pas rectangle en $A$.
Vérifions en $B$ : $\overrightarrow{BA}(-3;-1)$, $\overrightarrow{BC}(-2;2)$.
$\overrightarrow{BA}\cdot\overrightarrow{BC}=(-3)\times(-2)+(-1)\times2=6-2=4\neq0$.
Vérifions en $C$ : $\overrightarrow{CA}(-1;-3)$, $\overrightarrow{CB}(2;-2)$.
$\overrightarrow{CA}\cdot\overrightarrow{CB}=(-1)\times2+(-3)\times(-2)=-2+6=4\neq0$.
Le triangle n'est pas rectangle. (Erreur dans l'énoncé : il fallait $C(2;5)$ ? Vérifions : $A(1;2)$, $B(4;3)$, $C(2;5)$ donne $\overrightarrow{AB}(3;1)$, $\overrightarrow{AC}(1;3)$ non orthogonaux. En fait, $ABC$ n'est pas rectangle.)

\corrige{de l'exercice 7}
Un vecteur directeur de $(D)$ est $\vec{v}(4;3)$ (car $3x-4y+5=0$ donne $\vec{v}(4;3)$).
La droite cherchée a pour vecteur normal $\vec{v}$, donc une équation de la forme $4x+3y+c=0$.
Elle passe par $A(-2;3)$ : $4\times(-2)+3\times3+c=0\iff-8+9+c=0\iff c=-1$.
Donc l'équation est $4x+3y-1=0$.

\corrige{de l'exercice 8}
D'après Al-Kashi : $BC^2=AB^2+AC^2-2AB\cdot AC\cos\widehat{A}$.
$64=49+81-2\times7\times9\cos\widehat{A}\iff64=130-126\cos\widehat{A}\iff126\cos\widehat{A}=66\iff\cos\widehat{A}=\dfrac{66}{126}=\dfrac{11}{21}$.
$\widehat{A}=\arccos\left(\dfrac{11}{21}\right)\approx58^\circ$.

\corrige{de l'exercice 9}
$BC^2=AB^2+AC^2-2AB\cdot AC\cos\widehat{A}=25+36-2\times5\times6\times\cos120^\circ=61-60\times\left(-\dfrac12\right)=61+30=91$.
Donc $BC=\sqrt{91}$.
Aire : $\mathcal{A}=\dfrac12 AB\cdot AC\sin\widehat{A}=\dfrac12\times5\times6\times\sin120^\circ=15\times\dfrac{\sqrt3}{2}=\dfrac{15\sqrt3}{2}$.

\corrige{de l'exercice 10}
Soit $I$ le milieu de $[BC]$. $AI$ est la médiane issue de $A$.
D'après la formule de la médiane : $AB^2+AC^2=2AI^2+\dfrac{BC^2}{2}$.
$16+36=2AI^2+\dfrac{49}{2}\iff52=2AI^2+\dfrac{49}{2}\iff2AI^2=52-\dfrac{49}{2}=\dfrac{104-49}{2}=\dfrac{55}{2}\iff AI^2=\dfrac{55}{4}$.
Donc $AI=\dfrac{\sqrt{55}}{2}$.

\corrige{de l'exercice 11}
L'équation du cercle de centre $\Omega(2;-3)$ et de rayon $5$ est $(x-2)^2+(y+3)^2=25$.

\corrige{de l'exercice 12}
Le centre du cercle est le milieu $I$ de $[AB]$ : $I\left(\dfrac{1+5}{2};\dfrac{2+4}{2}\right